% theme: mechanical_energy
\MajorQuestion{力学的エネルギー}
\SetRowSpaceQanda

\Instruction{エネルギーの意味と位置エネルギーに関する次の問いに答えなさい。}
\QQ{ほかの物体に仕事をすることができる能力を\NamedBlank{energy}という。}{}
\QQ{\RefBlank{energy}の単位を何というか。}{}
\QQ{高い位置にある物体がもつ\RefBlank{energy}を\NamedBlank{potential_energy}という。}{}
\QQ{物体の高さを比べるとき、0として決める面を何というか。}{}
\QQ{物体にはたらく重力の大きさを何というか。}{}

\Instruction{位置エネルギーの大小と計算に関する次の問いに答えなさい。}
\QQ{同じ物体では、高さが高いほど\RefBlank{potential_energy}はどうなるか。}{}
\QQ{同じ高さでは、質量が大きいほど\RefBlank{potential_energy}はどうなるか。}{}
\QQ{重さ12Nの物体を0.75m高くすると、増える\RefBlank{potential_energy}は何Jか。}{}
\QQ{重さ4Nの物体が基準面から2.5mの高さにあるとき、\RefBlank{potential_energy}は何Jか。}{}
\QQ{重さ7.5Nの物体が0.8m低い位置に移動すると、減る\RefBlank{potential_energy}は何Jか。}{}

\Instruction{運動エネルギーに関する次の問いに答えなさい。}
\QQ{運動している物体がもつ\RefBlank{energy}を\NamedBlank{kinetic_energy}という。}{}
\QQ{\RefBlank{kinetic_energy}の大きさに関係する、物体の運動のはやさを表す量を何というか。}{}
\QQ{\RefBlank{kinetic_energy}の大きさに関係する、物体そのものの量を何というか。}{}
\QQ{同じ質量の小球では、速さが大きいほど\RefBlank{kinetic_energy}はどうなるか。}{}
\QQ{同じ速さの小球では、質量が大きいほど\RefBlank{kinetic_energy}はどうなるか。}{}

\Instruction{力学的エネルギーとその保存に関する次の問いに答えなさい。}
\QQ{\RefBlank{potential_energy}と\RefBlank{kinetic_energy}の和を\NamedBlank{mechanical_energy}という。}{}
\QQ{別の種類の\RefBlank{energy}に変わることを、\RefBlank{energy}の何というか。}{}
\QQ{摩擦力や空気の抵抗がないとき、\RefBlank{mechanical_energy}が一定に保たれることを何というか。}{}
\QQ{ジェットコースターが高い位置へ上がると、\RefBlank{potential_energy}はどうなるか。}{}
\QQ{摩擦力や空気の抵抗があると、\RefBlank{mechanical_energy}はどうなるか。}{}
